\section{Open questions}

The five open questions below emerged from a re-read of arXiv:1502.02677
after this replication was complete. Each is scoped so it could be
prosecuted as an independent follow-up study; each cites the specific
place in KLY that motivates the question and gives a concrete first
experimental step.

\begin{enumerate}
\item \textbf{Two-qubit gate extension.}\\
  \emph{Question.} Does RPE retain Heisenberg scaling for two-qubit
  entangling gate calibration (CNOT / CZ / cross-resonance) with
  realistic crosstalk, and how does the ladder generalize when the gate
  has multiple independent systematic error parameters (ZZ
  over-rotation, single-qubit phase leakage, control-target asymmetry)?\\
  \emph{Basis.} KLY Sec.~IV.C sketches nested single-qubit RPE for a
  universal gate set but does not directly address two-qubit gates. The
  paper's abstract restricts to a ``universal \emph{single-qubit}
  gate-set''; later work (Rudinger et al.~2017; Meier et al.~2019)
  extends to 2Q but was not touched in this replication.\\
  \emph{Next steps.} Extend \texttt{code/rpe\_sim.py} to a 2-qubit
  statevector simulator with a parameterized $\mathrm{CR}(\theta)$ plus
  residual $IX/IZ/ZI/ZZ$ terms; define cos/sin analogues via two-qubit
  Pauli measurements in a Bell basis; repeat the $K=1..14$ sweep and
  fit $\sigma$ vs $N$; compare against the tightened Rudinger~2017
  constants.

\item \textbf{Scalability under time-varying noise.}\\
  \emph{Question.} How does RPE precision degrade when the gate's
  over-rotation $\varepsilon$ drifts on the same timescale as ladder
  acquisition (e.g., slow $1/f$ flux noise on superconducting qubits)?\\
  \emph{Basis.} KLY assumes a fixed unknown phase $A$ during
  estimation; the Heisenberg-scaling proof implicitly requires
  stationarity across the $K$ generations. Drift is the dominant
  failure mode in real calibration workflows but is not modelled in the
  paper.\\
  \emph{Next steps.} Modify \texttt{rpe\_sim.py} to inject a stochastic
  $\varepsilon(t)$ with configurable OU-process or $1/f$ PSD; sweep the
  drift correlation time $\tau_{\text{drift}}$ against total ladder
  wall-clock $T_{\text{ladder}}$; identify the
  $\tau_{\text{drift}}/T_{\text{ladder}}$ ratio below which RPE loses
  its $N^{-1}$ advantage vs shot-noise-only; propose an
  interleaved-generation schedule that hedges against drift.

\item \textbf{Hybrid RPE + randomized benchmarking (RB).}\\
  \emph{Question.} Can a hybrid RPE+RB protocol combine RPE's sample
  efficiency for systematic parameters with RB's robustness to
  arbitrary stochastic noise, and what is the sample-complexity Pareto
  frontier?\\
  \emph{Basis.} KLY Sec.~VI explicitly positions RPE as an alternative
  to RB (Magesan et al.~2011); no head-to-head sample-complexity study
  was run in this replication.\\
  \emph{Next steps.} Add \texttt{code/rb\_baseline.py} running standard
  Clifford RB on the same $R_x(\pi/2+\varepsilon)$ circuit; measure the
  sample count needed for RB to detect the same $\varepsilon$ at
  3-$\sigma$ vs the RPE run at each $K$; plot RPE and RB on a common
  (samples, precision) axis; prototype a hybrid where RB bounds the
  stochastic error budget and RPE reads out the residual coherent
  phase.

\item \textbf{Hardware-noise-model sensitivity.}\\
  \emph{Question.} How sensitive is the RPE estimator to the specific
  hardware noise model ($T_1$ relaxation, $T_2$ dephasing, readout
  confusion, $|2\rangle$ leakage), and does the paper's SPAM-robustness
  proof (Sec.~IV) survive realistic superconducting/trapped-ion device
  parameters?\\
  \emph{Basis.} KLY Sec.~IV proves robustness to a specific
  additive-SPAM error model $(\delta_{\text{prep}},\delta_{\text{meas}})$;
  the paper does not simulate $T_1/T_2$/leakage/readout confusion. Our
  replication used a noiseless simulator (Sec.~4.1 of REPORT.md; C4
  ``Not tested''), so the SPAM claim is confirmed here only
  theoretically.\\
  \emph{Next steps.} Re-run the sweep against Qiskit \texttt{FakeManila}
  / \texttt{FakeGuadalupe} (real IBMQ calibration snapshots including
  readout confusion matrix); parametrize
  $T_1 \in [50\,\mu s, 200\,\mu s]$ and
  $T_2 \in [30\,\mu s, 150\,\mu s]$; identify the region where the RPE
  slope stays within 10\% of $-1.0$ vs where it degrades to the
  shot-noise plateau $-0.5$; compare against KLY's analytic
  $(\delta_{\text{prep}},\delta_{\text{meas}})$ bound.

\item \textbf{ML-decoder-boosted RPE.}\\
  \emph{Question.} Can a machine-learning (transformer / RNN) decoder
  trained on RPE bitstream patterns extract phase estimates with lower
  variance than the atan2 estimator (Eq.~V.3), particularly in the
  low-$M$ / low-shot regime where atan2 discretization noise dominates?\\
  \emph{Basis.} KLY uses a fixed
  $\operatorname{atan2}$-of-empirical-frequencies estimator, which is
  asymptotically optimal but not necessarily finite-shot optimal
  (Bayesian-MLE-for-phase-estimation work, e.g., Wiebe \& Granade 2016,
  suggests substantial gains at small $M$). No ML-decoder experiments
  were run.\\
  \emph{Next steps.} Generate a labeled training corpus by running
  \texttt{rpe\_sim.py} with known true $A$ across a dense grid; train
  (a) a Bayesian MLE estimator and (b) a small transformer taking the
  $(n_0^{(j)}, n_+^{(j)}, k_j)$ sequence as input and outputting
  $\hat A$; compare RMSE-vs-$N$ against the atan2 baseline; verify that
  the decoder respects Heisenberg scaling asymptotically while beating
  atan2 in the finite-shot regime.
\end{enumerate}
